% !TeX root = moduli-of-dynamics-on-projective-spaces.tex
\section{Moduli of dynamics on projective spaces}


The references of this section are \cite{Silverman-2012-ModuliSpacesArithmetic,Levy-2011-SpaceMorphismsProjective} and \Yang{}
For some basic knowledge of moduli problem and geometric invariant theory, I refer to \Yang{}

Throughout this section, let $\kk$ be an arbitrary field of characteristic zero and $\kkk$ be its algebraic closure without specific.




\subsection[status=draft,tag=0000000,label=\labelNFM{subsec:review-of-moduli-problem}]{Review of moduli problem and GIT}

Recall that a moduli problem is a functor.


Moduli space is a geometric space whose points represent isomorphism classes of certain geometric objects.
For example, fix an algebraically closed field \(\kkk\), all elliptic curves over \(\kkk\) can be classified by the \(j\)-invariant, which gives a bijection between isomorphism classes of elliptic curves and elements of \(\bbA^1(\kkk)\).
However, this classification is just a set-theoretic one.
We would like to have a geometric ``parameter space'' such that we can ``deform'' the objects continuously.
This is the initial motivation for moduli spaces.

In algebraic geometry, ``deforming objects continuously'' is usually described by flat families.
Hence the most perfect object to represent this moduli problem is such a flat family \(\calX \to M\), where \(M\) is a variety parameterizing all elliptic curves, and the fiber over each point \(m \in M(\kkk)\) is the elliptic curve corresponding to \(m\).
Furthermore, we hope that this family is universal, i.e., any other flat family \(\calY \to T\) of elliptic curves is obtained by pulling back \(\calX \to M\) along a unique morphism \(T \to M\).
(Despite in this example, such a perfect family does not exist; we will discuss this later.)
It can be described by the language of functors.
Consider the functor
\[ \calM: \Var_{\kkk}^{\op} \to \Set,\quad T \mapsto \{\text{flat families of elliptic curves over } T\}/\sim, \]
where \(\sim\) is the equivalence relation given by isomorphisms of families over \(T\).
If the family \(\calX \to M\) above exists, then it represents the functor \(\calM\).
In this case, we say that the moduli problem \(\calM\) is \emph{representable}, \(M\) is called a \emph{fine moduli space} and \(\calX \to M\) is called a \emph{universal object}.

Hence, to study a moduli problem, we follow the steps:
\begin{enumerate}[nosep]
	\item Define the moduli functor \(\calM\).
	\item Check whether \(\calM\) is representable.
\end{enumerate}

\begin{slogan}\label{slogan:moduli_problem_is_find_representable_functor}
	\hspace{1em}
	\begin{center}
		Moduli problem of parameterizing objects in $\Sigma$ \(\longleftrightarrow\) Find a functor $\calM$ with $\calM(\kkk) = \Sigma$ describing ``deformations'';

		Algebraic families over $T$ with fiber in $\Sigma$ $\longleftrightarrow$ elements in $\calM(T)$;

		Pullpack / base change along $f:T' \to T$ $\longleftrightarrow$ the map $\calM(f): \calM(T) \to \calM(T')$.
	\end{center}
	\Yang{}
\end{slogan}

\Yang{}
An important and we will use later example is the hilbert scheme and moduli space of endomorphisms. \Yang{}

Let us list some useful fine moduli spaces.

\begin{example}[label=\labelNFM{eg:hilbert-scheme}, status=draft, tag=0000000]
	Let $X$ be a projective variety over $\kk$ and $\calH$ an ample line bundle, $P$ is polynomial of $\bbQ$-coefficients.
	The \emph{Hilbert functor} is defined by
	\[ \calHilb_{(X,\calH,P)/\kk}(T) \coloneqq \left\{  \;\middle|\;  \right\} \]
	It is representable and represented by a scheme of finite type over $\kk$, called the \emph{Hilbert scheme}, denoted by $\Hilb_{(X,\calH,P)/\kk}$.
	\Yang{}
\end{example}

\begin{example}[label=\labelNFM{eg:moduli-functor-of-endmorphisms}, status=draft, tag=0000000]
	Let $X$ be a projective variety over $\kk$ and $\calH$ an ample line bundle.
	Its \emph{moduli functor of endomorphisms}
	\Yang{}
\end{example}

\begin{example}[label=\labelNFM{eg:moduli-of-functor-of-automorphisms}, status=draft, tag=0000000]
	Let $X$ be a projective variety over $\kk$ and $\calH$ an ample line bundle.
	Its \emph{moduli functor of numerically trivial automorphisms}
	\Yang{}
\end{example}





The main moduli functor we will care about is the moduli functor of dynamics on a variety.

\begin{definition}[label=\labelNFM{def:the-moduli-functor-of-dynamics-on-projective-spaces}, status=draft, tag=0000000]
	Let $\kk$ be a field of characteristic zero and $X$ a projective (and geometrically integral) variety over $\kk$.
	Let $\calPolDyn_{X,q}: \Sch_\kk^{\op} \to \Set$ be the functor defined by
	\[ T \mapsto \left\{ f:X\times T \to X \times T \;\middle|\; f_t: X \times \{t\} \to X \times \{t\} \text{ is } q\text{-polarized for every } t \in T \right\} / \sim. \]
	On this note, we will use $\calD_{n,d}$ to denote the moduli functor $\calPolDyn_{\bbP^n,d}$.
\end{definition}

However, for many questions we care about, the fine moduli space does exist.
An important statuation is that, if an object $X \in \calM(\kkk)$ has non-trivial automorphism group, then $\calM$ cannot have a fine moduli space.

\begin{proposition}[label=\labelNFM{prop:nontrivial-automorphism-group-means-no-fine-moduli-space}, status=draft, tag=0000000]
	Let $\calM: \Sch_\kk^\op \to \Set$ be a moduli functor.
	Suppose that there exists $X \in \calM(\kkk)$ such that its automorphism group in $\Sch_\kk$ is non-trivial, i.e., $\Aut_{\Sch_\kk}(X) \neq \{\id_X\}$.
	Then the fine moduli space of $\calM$ does not exit.
\end{proposition}
\begin{proof}
	\Yang{}
\end{proof}

Since the morphism $[-1]_E:E \to E$ is a non-trivial automorphism, above example of moduli space of elliptic curves has no fine moduli space; see also \cref{eg:affine-line-is-coarse-moduli-of-elliptic-curves}.

Hence, there is a weaken version, called \emph{coarse moduli space}, satisfying the following considition:
\begin{itemize}[nosep]
	\item there is an $1:1$ corresponding between the isomorphic classes parameterized by $\calM$ and the closed points in $M$;
	\item for every ``algebraic family $F \to T$ with fiber being objects parameterized by $\calM$'', there exists a unique morphism $f: T \to M$ induced by $F$ such that $F_t$ is isomorphic to the object parameterized by $f(t) \in M$.
\end{itemize}

In a rigorous language of functor, ``an algebraic family over $T$'' means an element $F$ in $\calM(T)$; a morphism corresponds to an element $f_F \in \Mor(T,M) = h_M(T)$.
Hence a coarse moduli space should consist of a scheme $M$ and a natural transformation $\eta: \calM \to h_M$ with $\eta(T):\calM(T) \to h_M(T), F \mapsto f_F$.
The the first condition above is just $\eta(\kkk): \calM(\kkk) \to h_M(\kkk)$ being bijective.

The remain issue is the uniqueness of $M$.
For example, if we have a morphism $g: N \to M$ which is bijective on the level of closed points, then $(N, g^* \circ \eta)$ is another pair satisfying the conditions.
(here $g^*:h_M \to h_N$ is the natural transformation given by composition by $g$.)
In a reasonable sense, since $N$ is onto $M$ by $g$, we regard that $M$ is the more ``suitable'' moduli space.
Hence we may require that the coarse moduli space $(M,\eta)$ with the universal property that other such pair $(N,\xi)$ should factor through it.
% On the other hand, note that such a factorization $\xi = \Phi \circ \eta$ must be given by a morphism $g:N \to M$ by Yoneda lemma.
% If require that such factorization is unique (as any other universal property), we naturally get the uniqueness of the morphism $f_F:T \to M$ above.
% For a general ``family'' $F \in \calM(T)$, we do need to require that all isomorphic classes appearing as fibers.
% Therefore, the universal property should also be stated without the assumption of $\xi(\kkk)$ being bijective.

In summary, we have the following definition.

\begin{definition}[label=\labelNFM{def:the-coarse-moduli-space}, status=draft, tag=0000000]
	Let $\calM$ be a moduli functor over $\kk$.
	A \emph{coarse moduli space} of $\calM$ is a pair $(M,\eta)$ with $M$ a $\kk$-scheme of finite type and $\eta: \calM \to h_M$ a natural transformation such that
	\begin{enumerate}
		\item for any other pair $(N,\xi \colon \calM \to h_{N})$, there exists a unique factorization
		      \[
			      \begin{tikzcd}
				      \calM & h_M \\
				      & h_{N}.
				      \arrow[from=1-1, to=1-2, "\eta"]
				      \arrow[from=1-2, to=2-2, "\exists !", gray]
				      \arrow[from=1-1, to=2-2, "\xi"']
			      \end{tikzcd}
		      \]
		\item $\eta(\kkk): \calM(\kkk) \to h_M(\kkk) = M(\kkk)$ is a bijection;
	\end{enumerate}
\end{definition}

By the universal property, if the coarse moduli space exists, then it is unique.

\begin{remark}[label=\labelNFM{rmk:fine-moduli-space-is-coarse-moduli-space-and-converse-if-universal-family-exists}, status=draft, tag=0000000]
	Let $\calM$ be a moduli functor.
	Suppose that $\calM$ has a fine moduli space, that is, $\calM$ is representable and represented by a scheme $M$ of finite type.
	Then easily see that $M$ is also a coarse moduli space of $\calM$.

	Conversely, suppose that $\calM$ has a coarse moduli space $(M,\eta)$.
	Then $M$ is a fine moduli space iff $\eta$ is a natural-isomorphism.
	This is also equivalent to that $\eta(T)$ is injective for all $T$ and there exists an element $\calX \in \calM(M)$ which gives the identity morphism $\eta(M)(\calX) = \id_M \in h_{M}(M)$.
\end{remark}

Many coarse moduli spaces are constructed by taking quotients of fine moduli spaces.
Let us introduce some notions of quotients in algebraic geometric.

\Yang{}
An important tool, and will be used later, to construct moduli space is the \emph{geometric invariant theory (GIT)}.

Here we provide a brief overview of geometric invariant theory (GIT).
For the propose of this notes, we only consider that the ground field $\kkk$ is algebraically closed and of characteristic zero, typically $\kkk = \bbC$.

\begin{definition}[Categorical and geometric quotients]
	Let $G$ be an algebraic group acting on a variety $X$ via $\sigma: G \times X \to X$.
	A \emph{categorical quotient} of $X$ by $G$ is a variety $Y$ together with a $G$-invariant morphism $\pi: X \to Y$ such that for any other variety $Z$ and any $G$-invariant morphism $f: X \to Z$, there exists a unique morphism $\tilde{f}: Y \to Z$ making the following diagram commute:
	\[
		\begin{tikzcd}
			G \acts X & Y \racts \id \\
			& Z \racts \id
			\arrow[from=1-1, to=1-2, "\pi"]
			\arrow[from=1-2, to=2-2, "\tilde{f}", gray]
			\arrow[from=1-1, to=2-2, "f"]
		\end{tikzcd}
	\]

	Given a categorical quotient $\pi: X \to Y$, we say that $Y$ is a \emph{geometric quotient} of $X$ by $G$ if the following conditions hold:
	\begin{enumerate}
		\item The variety $Y$ is equipped with the quotient topology induced by $\pi$, i.e., a subset $U \subseteq Y$ is open if and only if $\pi^{-1}(U)$ is open in $X$.
		\item The morphism $\pi$ is surjective, and the fibers of $\pi$ are precisely the $G$-orbits in $X$, i.e., for each $y \in Y$, we have $\pi^{-1}(y) = G \cdot x$ for some $x \in X$.
		\item The sheaf of regular functions on $Y$ is isomorphic to the sheaf of $G$-invariant regular functions on $X$, i.e., for any open subset $U \subseteq Y$, we have $\mathcal{O}_Y(U) \cong \mathcal{O}_X(\pi^{-1}(U))^G$.
	\end{enumerate}
	\Yang{}
\end{definition}
\Yang{}

In general, a categorical quotient may not be a geometric quotient, and the existence of quotients is not guaranteed.
\begin{example}
	Let $G = \bbG_m$ act on $X = \bbA^1$ by scaling, i.e., $t \cdot x = tx$ for $t \in \bbG_m$ and $x \in \bbA^1$.
	The categorical quotient of this action is a single point, since any $G$-invariant morphism from $\bbA^1$ to another variety must be constant.
	Easily, we can see that the categorical quotient is not a geometric quotient, since the fibers of the quotient map are not precisely the $G$-orbits in $X$.
\end{example}


Now let us restate the example of moduli functor of elliptic curves over $\kk$.

\begin{example}[label=\labelNFM{eg:affine-line-is-coarse-moduli-of-elliptic-curves}, status=draft, tag=0000000]
	Let $\calM$ be the functor given by
	\[ \calM: \Sch_{\kk}^{\op} \to \Set,\quad T \mapsto \left\{ (\pi:\calE \to T, e:T \to \calE \;\middle|\; \substack{\pi \text{ is flat and projective, } e \circ \pi = \id_T,\ \calE_t \text{ is an elliptic} \\ \text{ curve over $\rkk(t)$ with zero elements $e(t)$, } \forall t \in T}  \right\} \big/ \sim, \]

	We claim that $\bbA^1_\kkk$ is the coarse moduli space of $\calM$.
	\Yang{}
\end{example}







On GIT, we are particularly interested in the case where $G$ is a reductive group.

Recall that any affine linear algebraic group $G$ is a linear algebraic group.
For a linear algebraic group $G$ and $g \in G$, we say that $g$ is \emph{semisimple} if it is diagonalizable over the algebraic closure of the base field, and $g$ is \emph{unipotent} if all its eigenvalues are equal to 1.
We have the \emph{Jordan decomposition} of $g$ as $g = g_s g_u$, where $g_s$ is semisimple and $g_u$ is unipotent, and they commute with each other.

On the other hands, $G$ has a unique maximal normal solvable connected subgroup, which is called the \emph{radical} of $G$, denoted by $R(G)$.
All unipotent elements of $R(G)$ form a normal subgroup of $G$, called the \emph{unipotent radical} of $G$, denoted by $R_u(G)$.
The unipotent radical $R_u(G)$ can also be characterized as the unique maximal normal unipotent connected subgroup of $G$.
We say that $G$ is \emph{reductive} if its unipotent radical is trivial, i.e., $R_u(G) = \{e\}$, where $e$ is the identity element of $G$.
This is equivalent to saying that the radical $R(G)$ is a torus, i.e., $R(G) \cong (\bbG_m)^n$ for some $n \geq 0$.

\begin{example}
	The general linear group $\GL_n$ is reductive, since its radical is the center of $\GL_n$, which is isomorphic to $\bbG_m$.

	To see this, let $R = R(\GL_n)$ be the radical of $\GL_n$.
	By Lie-Kolchin theorem, $R$ is conjugate to a subgroup of the group $B$ of upper triangular matrices in $\GL_n$.
	Since $R$ is normal in $\GL_n$, it is contained in the intersection of all conjugates of $B$, which is the group of scalar matrices, i.e., $R \subseteq Z(\GL_n) \cong \bbG_m$.
	Note that we always have $Z(G)^\circ \subseteq R(G)$ for any linear algebraic group $G$, where $Z(G)^\circ$ is the connected component of the identity in the center of $G$.
	Hence $R$ is a torus, and thus $\GL_n$ is reductive.

	By similar argument, one can show that $R(\PGL_n) = \{e\}$, and thus $\PGL_n$ is also reductive.
\end{example}

For reductive groups, we have the following theorem to show its power.

\begin{theorem}[{{ref. \cite[Theorem 1.1 and Amplification 1.3]{Mumford-Fogarty-Kirwan-1994-GeometricInvariantTheorya}}}]
	Let $G$ be a reductive group acting on an affine variety $X$.
	Then the categorical quotient $X \to X / G$ exists and $X/G$ is also an affine variety over $\kk$.
	And it is a geometric quotient if and only if all $G$-orbits (of closed points) in $X$ are closed.
	\Yang{}
\end{theorem}

\begin{remark}[label=\labelNFM{rmk:quotient-by-nonreductive-group-may-be-not-of-finite-type}, status=draft, tag=0000000]
	Let $G$ be a general algebraic group acting on an affine scheme $X = \Spec A$.
	$A^G$ may be not of finite type over $\kk$.
	\Yang{}
\end{remark}

\begin{remark}[label=\labelNFM{rmk:quotient-for-projective-case-GIT}, status=draft, tag=0000000]
	For projective case, there is also a whole theory to deal the quotient by reductive group, namely, the \emph{geometric invariant theory}.
	The standard reference is \cite{Mumford-Fogarty-Kirwan-1994-GeometricInvariantTheorya}.
\end{remark}



% For projective cases, the case is a bit more complicated, and we need to introduce the notion of $G$-linearization and (semi)stable points.
%
% \begin{definition}
% 	Let $G$ be an algebraic group acting on a variety $X$ and let $\calL$ be a line bundle on $X$.
% 	A \emph{$G$-linearization} of $\calL$ is an isomorphism $\phi: \sigma^* \calL \to p_2^* \calL$ of line bundles on $G \times X$, where $\sigma: G \times X \to X$ is the action map and $p_2: G \times X \to X$ is the projection onto the second factor.
% \end{definition}
%
% Inspiringly, assume that $\calL$ is very ample with an embedding $\iota: X \hookrightarrow \bbP^n$.
% Then given a $G$-linearization $\phi: \sigma^* \calL \to p_2^* \calL$, by restricting $\phi$ to the $\{g\} \times X$ for each $g \in G$, we get an isomorphism $\phi_g: \sigma_g^* \calL \to \calL$, where $\sigma_g: X \to X$ is the action of $g$ on $X$.
% Hence we induce a (contravariant) linear action of $G$ on $H^0(X, \calL)$, and thus a linear action of $G$ on $\bbP^n$.
% In this case, a $G$-linearization of $\calL$ will induce a $G$-equivariant embedding $\iota: X \hookrightarrow \bbP^n$.
%
% \begin{definition}
% 	Let $G$ be a reductive group acting on a projective variety $X$ and let $\calL$ be a $G$-linearized line bundle on $X$.
% 	A point $x \in X$ is called \emph{semistable} if there exists a $G$-invariant section $s \in H^0(X, \calL^{\otimes n})^G$ for some $n > 0$ such that $s(x) \neq 0$ and $X_s := \{y \in X \mid s(y) \neq 0\}$ is affine.
% 	A semistable point $x \in X$ is called \emph{stable} if the action of $G$ on $X_s$ is closed, i.e., the orbit $G \cdot y$ is closed in $X_s$ for all $y \in X_s$.
% \end{definition}
%
% The set of semistable points and stable points are open in $X(\kkk)$, and we will also denote the corresponding open subvarieties by $X^{\text{ss}}(\calL)$ and $X^{\text{s}}(\calL)$ respectively.
%
% Then we have the following theorem in projective cases.
%
% \begin{theorem}[{{ref. \cite[Theorem 1.10 and comments under Amplification 1.11]{Mumford-Fogarty-Kirwan-1994-GeometricInvariantTheorya}}}]
% 	Let $G$ be a reductive group acting on a projective variety $X$ and let $\calL$ be a $G$-linearized line bundle on $X$.
% 	\begin{enumerate}
% 		\item there exists a categorical quotient $\pi: X^{\text{ss}}(\calL) \to X^{\text{ss}}(\calL) / G$;
% 		\item $X^{\text{ss}}(\calL) / G \cong \Proj \left( \bigoplus_{n \geq 0} H^0(X, \calL^{\otimes n})^G \right)$;
% 		\item the restriction of $\pi$ to $X^{\text{s}}$ is a geometric quotient $\pi: X^{\text{s}}(\calL) \to X^{\text{s}}(\calL) / G$.
% 	\end{enumerate}
% \end{theorem}
%
% The next problem is to determine the (semi)stable points.
%
% \begin{construction}
% 	Settings as the theorem above.
% 	Let $l: \bbG_m \to G$ be a homomorphism of algebraic groups, which is called a \emph{$1$-parameter subgroup} of $G$.
% 	Then we have an action of $\bbG_m$ on $X$ via $l$, and thus we can define the limit point $\lim_{t \to 0} l(t) \cdot x$ for any $x \in X$.
% 	More exactly, given an $1$-parameter subgroup $l: \bbG_m \to G$, we can define a morphism $\rho_{x, l}: \bbG_m \to X$ by $\rho_{x, l}(t) = l(t) \cdot x$.
% 	And then by the valuative criterion of properness, we can uniquely extend $\rho_{x, l}$ to a morphism $\tilde{\rho}_{x, l}: \bbA^1 \to X$, and thus we can define the limit point $\lim_{t \to 0} l(t) \cdot x = \tilde{\rho}_{x, l}(0)$.
% 	This limit point will be denoted by $l(0) \cdot x$.
%
% 	For every $\alpha \in \bbG_m$, we can shift the $1$-parameter subgroup $l$ by $\alpha$ to get another $1$-parameter subgroup $l_\alpha: \bbG_m \to G$ defined by $l_\alpha(t) = l(\alpha t)$.
% 	Note that $l_\alpha(0) \cdot x = l(0) \cdot x$ for all $\alpha \in \bbG_m$ since $\tilde{\rho}_{x, l_\alpha}(t) = \tilde{\rho}_{x, l}(\alpha t)$.
% 	Hence $l(\alpha) \cdot (l(0) \cdot x) = l(0) \cdot x$ for all $\alpha \in \bbG_m$.
%
% 	Note that since $\calL$ is $G$-linearized, we have an induced action of $G$ on the vector bundle $\bbV(\calL)$, and thus we have an action of $\bbG_m$ on the fiber $\calL_{l(0) \cdot x}$.
% 	This gives a character $\chi: \bbG_m \to \bbG_m$ defined by $l(\alpha) \cdot v = \chi(\alpha) v$ for all $v \in \bbV(\calL)|_{l(0)\cdot x}$ and $\alpha \in \bbG_m$.
% 	Identifying $\Hom(\bbG_m, \bbG_m) \cong \bbZ$, we can write $\chi(\alpha) = \alpha^r$ for some $r \in \bbZ$.
% 	Then we set $\mu^\calL(x, l) = r$.
% \end{construction}
%
% \begin{remark}
% 	Note that $\bbV(\calL) = \Spec \Sym^\bullet \calL$ is in the sense of Grothendieck, and the fiber $\bbV(\calL)|_x$ is dual to the fiber $\calL|_x$.
% 	Hence our definition of $\mu^\calL(x, l)$ coincides with the one in \cite[Definition 2.2]{Mumford-Fogarty-Kirwan-1994-GeometricInvariantTheorya} although we have a sign difference.
% 	\Yang{}
% \end{remark}
%
% \begin{theorem}[{{ref. \cite[Theorem 2.1]{Mumford-Fogarty-Kirwan-1994-GeometricInvariantTheorya}}}]
% 	Let $G$ be a reductive group acting on a projective variety $X$ and let $\calL$ be a $G$-linearized line bundle on $X$.
% 	Then we have the following criterion for (semi)stability:
% 	\begin{enumerate}
% 		\item A point $x \in X$ is semistable if and only if $\mu^\calL(x, l) \geq 0$ for all $1$-parameter subgroups $l: \bbG_m \to G$.
% 		\item A point $x \in X$ is stable if and only if $\mu^\calL(x, l) > 0$ for all $1$-parameter subgroups $l: \bbG_m \to G$.
% 	\end{enumerate}
% \end{theorem}
%
% If $X = \bbP^n$ and $\calL = \calO_{\bbP^n}(1)$, then we can give a more explicit description of $\mu^\calL(x, l)$.
%
% \begin{proposition}[{{ref. \cite[Example 2.3]{Mumford-Fogarty-Kirwan-1994-GeometricInvariantTheorya}}}]
% 	Let $G$ be a reductive group act on $X = \bbP^n$ and $\calL = \calO_{\bbP^n}(1)$.
% 	Then for any $1$-parameter subgroup $l: \bbG_m \to G$, there exists a coordinate system such that $l(t) = \diag(t^{r_0}, t^{r_1}, \ldots, t^{r_n})$ with $\sum_{i=0}^n r_i = 0$.
% 	Under this coordinate system, for any point $x = [x_0: x_1: \ldots: x_n] \in \bbP^n$, we have
% 	\[
% 		\mu^\calL(x, l) = -\min\{r_i \mid x_i \neq 0\}.
% 	\]
% \end{proposition}
%







\subsection[status=draft,tag=0000000,label=\labelNFM{subsec:automorphism-group-and-construction-of-coarse-moduli}]{Automorphism of group and construction of coarse moduli}


% \subsection[status=draft,tag=0000000,label=\labelNFM{subsec:automorphism-group-of-dynamics}]{Automorphism group of dynamics on projective spaces}


First we show that the automorphism group of a dynamics $f: \bbP^n \to \bbP^n$ is finite.

\begin{definition}[label=\labelNFM{def:automorphism-group-of-dynamics}, status=draft, tag=0000000]
	Let $f: \bbP^n \to \bbP^n$ a surjective endomorphism over a field $\kk$.
	Then its automorphism group, is defined as an algebraic group by
	\[ \Aut_{f/\kk}(\kkk) \coloneqq \left\{ g \in \Aut_{\bbP^n/\kk} = \PGL_n() \;\middle|\; f \circ g = g \circ f \right\}.  \]
\end{definition}

\begin{remark}[label=\labelNFM{rmk:scheme-theoretic-automorphism-group-and-categroy-theoretic}, status=draft, tag=0000000]
	Note that for , we consider its automorphism group as a \emph{algebraic group}.
	And ...
\end{remark}




\begin{proposition}[label=\labelNFM{prop:automorphism-group-of-dynamics-on-pn-is-finite}, status=draft, tag=0000000]
	Let $f: \bbP^n \to \bbP^n$ a surjective endomorphism over a field $\kk$.
	Then its automorphism is finite, that is,
	\[ \# \Aut_{f/\kk}(\kkk) = \# \left\{ g \in \Aut(\bbP^n/\kk) = \PGL_n(\kk) \;\middle|\; f \circ g = g \circ f \right\} < \infty.  \]
\end{proposition}
\begin{proof}
	\Yang{}
\end{proof}

\begin{example}[label=\labelNFM{eg:power-map-has-nontrivial-automorphism-by-unit-roots}, status=draft, tag=0000000]
	Consider the power map $f: z \mapsto z^d$ for $d \geq 3$.
	Then $\mu_{d-1} \in \Aut_{f/\kk}(\kkk)$ since $m_\xi \circ f \circ m_{\xi^{-1}} (z) = \xi \cdot \xi^{-d} z^{d} = z^d = f(z)$ for $\xi \in \mu_{d-1}$ and $m_a$ is the map given by muliplication of $a$.
\end{example}

\begin{example}[label=\labelNFM{eg:power-map-has-nontrivial-automorphism-by-permutation}, status=draft, tag=0000000]
	More generally, let $f: [x_0:\cdots:x_n] \mapsto [x_0^d:\cdots:x_n^d]$.
	Then any involution $\tau \in \frakS_{n+1}$, i.e., an permutation of $n+1$ symbol satisfying $\tau^2 = \id$, induces an automorphism $g_\tau \in \Aut_{\bbP^n}(\kk)$ by $[x_0:\cdots : x_n] \mapsto [x_{\tau(0)}:\cdots:x_{\tau(n)}]$ which is commutative with $f$.
	In particular, $f$ has non-trivial automorphism group and $\calD_{d,n}$ can't have a fine moduli space.
\end{example}

Moreover, the bound is uniform in term of dimension $n$ and degree $d$.

\begin{theorem}[label=\labelNFM{thm:uniform-bound-for-the-auto-group-of-dynamics-on-pn}, status=draft, tag=0000000]
	There exists a number $C = C_{n,d}$ depending only on $n$ and $d$ such that for any field $\kk$ and $[f] \in \calD_{n,d}(\kk)$,
	\[ \# \Aut([f]/\kk) \leq C. \]
\end{theorem}
\begin{proof}
	\Yang{}
\end{proof}


Now let us turn to construct the coarse moduli of dynamics on $\bbP^n$ by GIT.
The \cref{prop:automorphism-group-of-dynamics-on-pn-is-finite} is enough to show the existence of the coarse moduli space.

\begin{theorem}[label=\labelNFM{thm:existence-of-coarse-moduli-of-dynamics-on-projective-space}, status=draft, tag=0000000]
	The coarse moduli space $D_{n,d}$ of the moduli functor $\calD_{n,d}$ exists and is an affine variety of dimension $(n+1)C_{n+d}^d - (n+1)^2$.
\end{theorem}
\begin{proof}
	\Yang{}
\end{proof}

A good example is the degree two dynamics on projective line.

\begin{example}[label=\labelNFM{eg:moduli-of-degree-two-maps-on-p1}, status=draft, tag=0000000]
	Let us consider $D_{1,2}$, i.e., the coarse moduli space of degree two endomorphism on $\bbP^1$.
	We claim that $D_{1,2} \cong \bbA^2$.


	\Yang{}
\end{example}

\begin{remark}[label=\labelNFM{rmk:rational-and-unirational-of-coarse-moduli-of-dynamics-on-pn}, status=draft, tag=0000000]
	The coarse moduli spaces $\calD_{n,d}$ is unirational since \Yang{}
	Moreover, \Yang{}
\end{remark}



% \subsection[status=draft,tag=0000000,label=\labelNFM{subsec:construction-of-coarse-moduli-of-dynamics-on-projective-space}]{Construction of coarse moduli}


% \begin{theorem}[label=\labelNFM{thm:coarse-moduli-of-dynamics-on-projective-line-is-affine-and-rational}, status=draft, tag=0000000]
% 	Let $\kk$ be a field and $D_{n,d}$ be the coarse moduli space of $d$-polarized endomorphism on projective space $\bbP^n$.
% 	Then
% 	\begin{enumerate}
% 		\item The variety $D_{n,d}$ is an affine variety of dimension $(n+1)C_{n+d}^d - (n+1)^2$.
% 		\item The variety $D_{1,d}$ is rational for all $d$.
% 	\end{enumerate}
% \end{theorem}
% \begin{proof}
% 	\Yang{}
% \end{proof}




% \subsection[status=draft,tag=0000000,label=\labelNFM{subsec:focus-on-projective-lines}]{Focus on projective line}

% \begin{definition}[label=\labelNFM{def:pcf-map-on-pn}, status=draft, tag=0000000]
% 	Let $f: \bbP^n \to \bbP^n$ be a \Yang{}
% 	We say that $f$ is PCF if \Yang{}
% \end{definition}
%
% \begin{theorem}[label=\labelNFM{thm:pcf-points-defined-over-Qbar}, status=draft, tag=0000000]
% 	Let $\kk$ be any field over $\bbQ$ and $[f] \in \calP_{1,d}(\kk)$ be a PCF point.
% 	Then $[f]$ is indeed defined over $\bar{\bbQ}$, i.e., there exists that $[f_{\bar{\bbQ}}] \in \calP_{1,d}(\bar{\bbQ})$ such that $f = f_{\bar{\bbQ}} \times_{\bar{\bbQ}} \kk$.
% 	\Yang{}
% \end{theorem}
%
%




